\section{Sistem Bilangan dan Konversi}
\label{sec:sistem-bilangan}

Komputer menyimpan seluruh data sebagai deretan bit. Angka yang
ditulis manusia dalam basis sepuluh harus diubah terlebih dahulu
menjadi bentuk yang dapat disimpan dan diolah prosesor. Karena itu,
pemahaman sistem bilangan bukan pelengkap, melainkan syarat untuk
membaca isi register, memeriksa memori, dan menulis data pada
program assembly.

Bagian ini diringkas dari \cite{sto2001}, BAB I, hlm.~6--8, dan
BAB V, hlm.~23--24, serta \cite{gunawan2019}, Bagian 14.1,
hlm.~127--128.

\subsection{Notasi Posisional}

Pada notasi posisional, nilai suatu angka ditentukan oleh angka itu
sendiri dan posisinya. Untuk bilangan dengan basis $b$, angka pada
posisi ke-$i$ (dihitung dari kanan, mulai nol) memiliki bobot $b^i$.
Secara umum, bilangan $d_n d_{n-1} \ldots d_1 d_0$ dalam basis $b$
bernilai

\[
    d_n b^n + d_{n-1} b^{n-1} + \cdots + d_1 b^1 + d_0 b^0.
\]

Contohnya, bilangan desimal $123_{10}$ bernilai
$1 \times 10^2 + 2 \times 10^1 + 3 \times 10^0 = 123$. Aturan yang
sama berlaku untuk basis lain; yang berubah hanya nilai basisnya.

\subsection{Empat Basis yang Dipakai}

Bahasa assembly 8086 mengenal empat jenis bilangan: biner, oktal,
desimal, dan heksadesimal. Setiap basis memiliki jumlah digit yang
berbeda.

\begin{itemize}
    \item \term{Biner} (basis 2) memakai digit 0 dan 1.
    \item \term{Oktal} (basis 8) memakai digit 0 sampai 7.
    \item \term{Desimal} (basis 10) memakai digit 0 sampai 9.
    \item \term{Heksadesimal} (basis 16) memakai digit 0 sampai 9
          ditambah A, B, C, D, E, dan F untuk nilai 10 sampai 15.
\end{itemize}

Kesetaraan digit pada keempat basis ditunjukkan pada
Tabel~\ref{tab:digit-bilangan}.

\begin{table}[H]
    \centering
    \caption{Kesetaraan digit pada bilangan biner, oktal, desimal,
    dan heksadesimal.}
    \label{tab:digit-bilangan}
    \small
    \begin{tabular}{@{}cccc@{}}
        \toprule
        \textbf{Desimal} & \textbf{Biner (4 bit)} &
        \textbf{Oktal} & \textbf{Heksadesimal} \\
        \midrule
        0 & \bin{0000} & 0 & 0 \\
        1 & \bin{0001} & 1 & 1 \\
        2 & \bin{0010} & 2 & 2 \\
        3 & \bin{0011} & 3 & 3 \\
        4 & \bin{0100} & 4 & 4 \\
        5 & \bin{0101} & 5 & 5 \\
        6 & \bin{0110} & 6 & 6 \\
        7 & \bin{0111} & 7 & 7 \\
        8 & \bin{1000} & 10 & 8 \\
        9 & \bin{1001} & 11 & 9 \\
        10 & \bin{1010} & 12 & A \\
        11 & \bin{1011} & 13 & B \\
        12 & \bin{1100} & 14 & C \\
        13 & \bin{1101} & 15 & D \\
        14 & \bin{1110} & 16 & E \\
        15 & \bin{1111} & 17 & F \\
        \bottomrule
    \end{tabular}
\end{table}

\subsection{Konversi Antarjenis Bilangan}

Konversi dapat dikerjakan dengan menghitung bobot posisi atau dengan
mengelompokkan bit. Pilihan cara bergantung pada arah konversi.

\subsubsection{Biner ke Desimal}

Kalikan setiap bit dengan bobotnya, lalu jumlahkan. Contoh untuk
\bin{11110011}:

\[
    1 \times 2^7 + 1 \times 2^6 + 1 \times 2^5 + 1 \times 2^4
    + 0 \times 2^3 + 0 \times 2^2 + 1 \times 2^1 + 1 \times 2^0
    = 128 + 64 + 32 + 16 + 2 + 1 = 243.
\]

\subsubsection{Desimal ke Biner}

Bagilah bilangan dengan 2 secara berulang, catat sisanya, lalu baca
sisa dari bawah ke atas. Contoh untuk 243:

\begin{center}
\begin{tabular}{@{}ccc@{}}
    \toprule
    \textbf{Pembagian} & \textbf{Hasil} & \textbf{Sisa} \\
    \midrule
    $243 \div 2$ & 121 & 1 \\
    $121 \div 2$ & 60 & 1 \\
    $60 \div 2$ & 30 & 0 \\
    $30 \div 2$ & 15 & 0 \\
    $15 \div 2$ & 7 & 1 \\
    $7 \div 2$ & 3 & 1 \\
    $3 \div 2$ & 1 & 1 \\
    $1 \div 2$ & 0 & 1 \\
    \bottomrule
\end{tabular}
\end{center}

Sisa dibaca dari bawah ke atas sehingga hasilnya \bin{11110011}.

\subsubsection{Biner ke Oktal dan Heksadesimal}

Kelompokkan bit dari kanan. Untuk oktal, setiap kelompok berisi tiga
bit; untuk heksadesimal, setiap kelompok berisi empat bit. Contoh
\bin{11110011}:

\begin{itemize}
    \item Oktal: \bin{011} \bin{110} \bin{011} $\rightarrow$
          3, 6, 3, sehingga hasilnya $363_8$.
    \item Heksadesimal: \bin{1111} \bin{0011} $\rightarrow$
          F, 3, sehingga hasilnya \hex{F3}.
\end{itemize}

\subsubsection{Oktal atau Heksadesimal ke Desimal}

Gunakan bobot posisi dengan basis 8 atau 16. Contoh untuk $363_8$:

\[
    3 \times 8^2 + 6 \times 8^1 + 3 \times 8^0
    = 192 + 48 + 3 = 243.
\]

Contoh untuk \hex{F3}:

\[
    15 \times 16^1 + 3 \times 16^0 = 240 + 3 = 243.
\]

\subsubsection{Desimal ke Heksadesimal}

Bagilah bilangan dengan 16 secara berulang, catat sisanya, lalu
terjemahkan sisa 10--15 menjadi A--F. Untuk 243:

\begin{center}
\begin{tabular}{@{}ccc@{}}
    \toprule
    \textbf{Pembagian} & \textbf{Hasil} & \textbf{Sisa} \\
    \midrule
    $243 \div 16$ & 15 & 3 \\
    $15 \div 16$ & 0 & 15 (F) \\
    \bottomrule
\end{tabular}
\end{center}

Sisa dibaca dari bawah ke atas sehingga hasilnya \hex{F3}. Ringkasan
keempat representasi untuk nilai 243 ditunjukkan pada
Tabel~\ref{tab:contoh-basis}.

\begin{table}[H]
    \centering
    \caption{Representasi nilai 243 pada empat basis.}
    \label{tab:contoh-basis}
    \small
    \begin{tabular}{@{}lll@{}}
        \toprule
        \textbf{Basis} & \textbf{Penulisan} & \textbf{Keterangan} \\
        \midrule
        Biner & \bin{11110011} & delapan bit \\
        Oktal & $363_8$ & tiga digit oktal \\
        Desimal & 243 & tanpa akhiran \\
        Heksadesimal & \hex{F3} & dua digit heksadesimal \\
        \bottomrule
    \end{tabular}
\end{table}

\subsection{Penulisan Literal Bilangan pada Assembler}

Assembler perlu mengetahui basis sebuah bilangan. Menurut
\cite{sto2001}, BAB V, hlm.~23--24, penulisan yang dipakai adalah
sebagai berikut.

\begin{itemize}
    \item \textbf{Desimal:} boleh tanpa akhiran atau diberi akhiran
          \code{D}, misalnya \code{243} atau \code{243D}.
    \item \textbf{Biner:} diberi akhiran \code{B}, misalnya
          \code{11110011B}.
    \item \textbf{Heksadesimal:} diberi akhiran \code{H}. Bila digit
          pertama berupa huruf A--F, tambahkan angka 0 di depan agar
          assembler tidak membacanya sebagai label, misalnya
          \code{0F3H}.
    \item \textbf{Karakter:} ditulis di antara tanda petik, misalnya
          \code{'A'} atau \code{"A"}.
\end{itemize}

Penulisan oktal tidak dipakai pada contoh program bab ini karena
assembler yang dipakai dalam buku lebih sering memakai bentuk biner,
desimal, dan heksadesimal. Nilai oktal tetap dilatih melalui konversi
dan dapat dituliskan kembali dalam bentuk desimal atau heksadesimal
sebelum dimasukkan ke program.

\begin{peringatanbox}
\vspace{1.1em}
Tiga kesalahan yang sering terjadi pada penulisan bilangan:
\begin{enumerate}
    \item Menulis \code{F3H} tanpa nol di depan. Assembler dapat
          menganggapnya sebagai nama label, bukan bilangan.
    \item Menulis angka 2--9 pada bilangan biner, misalnya
          \code{1012B}. Bilangan biner hanya boleh memuat 0 dan 1.
    \item Menganggap \code{10} dan \code{10H} bernilai sama. Yang
          pertama adalah sepuluh desimal, yang kedua adalah enam belas
          desimal.
\end{enumerate}
\end{peringatanbox}
